\documentclass{article}

% Language setting
% Replace `english' with e.g. `spanish' to change the document language
\usepackage[czech]{babel}

% Set page size and margins
% Replace `letterpaper' with `a4paper' for UK/EU standard size
\usepackage[a4paper,top=2cm,bottom=2cm,left=3cm,right=3cm,marginparwidth=1.75cm]{geometry}

\begin{document}
\hfill Jan Romanovský
\begin{center}
    \textbf{Kvantová podivnost}
\end{center}

Na začátku je krátce zopakován obsah předchozí přednášky, která byla první v cyklu Kvantový svět -- úvod to kvantové problematiky a pojmy maximální sady pozorovatelných, principu neurčito\-sti, představa možných stavů jako báze prostoru a s tím spojené představy meření jako projekce na podprostor, lineárního skládání stavů, určení fáze stavu.

Následuje popis normalizace stavu a vysvětlení rozdílu mezi superpozicí a směsí stavů.

Dále se vše převede do matematické řeči skalárního součinu reprezentací stavů, řeší se měření jako operátory na Hilbertově prostoru, fyzikální význam různých vlastností této matematické reprezentace (vlastní čísla, vlastní vektory, komutativita násobení, střední hodnota).

Pomocí této reprezentace se pak ptáme na časový vývoj systému, když ho neměříme -- zjistíme, že je to \uv{unitární} vývoj, který je deterministický.

Poté se dostáváme do druhé poloviny přednášky, kde se řeší problém zpožděného měření, nejdřív teoreticky a pak i s popisem, jak je zpožděné měření prováděno prakticky, obohaceno o vsuvku o J. A. Wheelerovi i s osobním setkáním prof. Krtouše.

Koncept zpožděného měření je dále využit v myšlenkovém experimentu neinterakčního měření bomby a rozvinut možností smazání výsledku měření, což se ukáže jako konzistentní s předchozími poznatky.

Nakonec je zmíněna slavná Schrödingerova kočka a její variace s Wignerovou přítelkyní, které řeší, kdy dochází ke kolapsu superpozice, a co se děje, když je víc pozorovatelů.

\end{document}
